\documentclass[11pt]{article}

\usepackage[margin=2.5cm]{geometry}
\usepackage{amsmath, amssymb, amsthm, mathtools}
\usepackage{enumitem}
\usepackage{hyperref}

\newtheorem{theorem}{Theorem}[section]
\newtheorem{proposition}[theorem]{Proposition}
\newtheorem{lemma}[theorem]{Lemma}
\newtheorem{corollary}[theorem]{Corollary}
\theoremstyle{definition}
\newtheorem{remark}[theorem]{Remark}
\newtheorem{conjecture}[theorem]{Conjecture}
\newtheorem{hypothesis}[theorem]{Hypothesis}

\newcommand{\One}{\textsc{One}}
\newcommand{\All}{\textsc{All}}
\newcommand{\N}{\mathbb{N}}
\newcommand{\R}{\mathbb{R}}
\newcommand{\Z}{\mathbb{Z}}

\title{Towards infinitely many local extrema of the optimal\\
  value $w_{n,p}$ of the all-heads coin game for $p<\tfrac12$}
\author{Working note --- P.\ Pfaffelhuber / Claude}
\date{\today}

\begin{document}
\maketitle

\begin{abstract}
  This note concerns the conjecture that, for $p<\tfrac12$, the optimal
  winning probability $w_{n,p}$ of the all-heads coin game has
  infinitely many strict local maxima and minima as a function
  of~$n$. We reduce the conjecture to two one-sided statements.
  The first --- \emph{$w_{n,p}$ is not eventually non-decreasing} ---
  is proved here unconditionally and elementarily. The second ---
  \emph{$w_{n,p}$ is not eventually non-increasing} --- is reduced, via
  an \emph{exact} transfer lemma, to the log-periodic oscillation of
  the strategy-\All{} value $b_{n,p}$ that is derived (modulo two
  standard analytic gaps) in the companion note
  \texttt{strategy\_all\_asymptotics.tex}. We record carefully which
  steps are rigorous and which remain open, and isolate the single new
  analytic point that a complete proof would still require.
\end{abstract}

\section{Setup, standing facts, and the conjecture}\label{sec:setup}

Fix $p\in(0,1)$, write $q:=1-p$, and let $w_n:=w_{n,p}$ denote the
optimal winning probability of the all-heads coin game with $n$ coins.
By definition $w_0=1$, and for $n\ge 1$ the Bellman equation reads
\begin{equation}\label{eq:bellman}
  w_n \;=\; p^n + \sum_{j=1}^{n-1}\binom{n}{j}p^{\,n-j}q^{\,j}\,
  \max_{j\le m\le n-1} w_m ,
\end{equation}
the empty sum at $n=1$ giving $w_1=p$.

\paragraph{Standing facts.} We use the following throughout, all
elementary or established in the main manuscript.
\begin{enumerate}[label=(F\arabic*),leftmargin=3.2em]
\item\label{f:bounds} $0<w_n\le 1$ for every $n$ (it is a probability,
  and $w_n\ge p^n>0$).
\item\label{f:pq} For $p<\tfrac12$ one has $q>\tfrac12>p$; hence
  $q>p$, $p/q<1$, and $2p<1$.
\item\label{f:binom} $\displaystyle\sum_{j=0}^{n}\binom nj p^{n-j}q^j=1$,
  $\displaystyle\sum_{j=1}^{n-1}\binom nj p^{n-j}q^j=1-p^n-q^n$, and
  $\displaystyle\sum_{j=0}^{n-1}\binom nj p^{n-j}q^j=1-q^n$.
\item\label{f:strat} The strategy-\One{} and strategy-\All{} values
  satisfy the \emph{linear} recursions
  \[
    a_n = p^n+(1-p^n-q^n)\,a_{n-1},\qquad
    b_n = \sum_{j=0}^{n-1}\binom nj p^{n-j}q^j\,b_j ,
  \]
  with $a_0=b_0=1$, and $w_n\ge\max(a_n,b_n)$ for all $n$ (the optimum
  dominates any fixed strategy).
\end{enumerate}

\paragraph{The conjecture.} For $p>\tfrac12$ the sequence $n\mapsto
w_{n,p}$ is strictly increasing (Theorem~3.1 of the main paper), and
for $p=\tfrac12$ it is constant. The case $p<\tfrac12$ is the
interesting one.

\begin{conjecture}\label{conj:main}
  For every $p\in(0,\tfrac12)$ the sequence $n\mapsto w_{n,p}$ is
  \emph{not eventually monotone}; equivalently, it has infinitely many
  strict local maxima and infinitely many strict local minima.
\end{conjecture}

The two formulations agree up to the (expected, and in all computed
data confirmed) genericity assumption that $w_{n,p}\ne w_{n+1,p}$ for
all~$n$; see Lemma~\ref{lem:reduction} and
Remark~\ref{rem:strictness}. Conjecture~\ref{conj:main} is consistent
with the numerical extremum data of the main manuscript (for $p=0.45$,
local minima at $n=6,17$ and a local maximum at $n=15$; for $p=0.42$,
minima at $n=7,13$ and a maximum at $n=9$) and with van~Doorn's
description of the optimal strategy for $p<\tfrac12$ as ``infinitely
complex''.

\section{Reduction to two one-sided statements}\label{sec:reduction}

For a real sequence $(x_n)$ put $d_n:=x_{n+1}-x_n$. Call $n\ge 1$ a
\emph{local maximum} if $x_{n-1}\le x_n\ge x_{n+1}$ and a \emph{strict}
local maximum if both inequalities are strict; symmetrically for
minima. The sequence is \emph{eventually non-decreasing} if $d_n\ge 0$
for all large~$n$, \emph{eventually non-increasing} if $d_n\le 0$ for
all large~$n$, and \emph{eventually monotone} if it is one of the two.

\begin{lemma}\label{lem:reduction}
  For any sequence $(x_n)$ the following are equivalent.
  \begin{enumerate}[label=\rm(\roman*),leftmargin=2.4em]
  \item $(x_n)$ is not eventually monotone.
  \item $d_n>0$ for infinitely many $n$ \emph{and} $d_n<0$ for
    infinitely many $n$.
  \item $(x_n)$ is neither eventually non-decreasing nor eventually
    non-increasing.
  \end{enumerate}
  Any of these implies that $(x_n)$ has infinitely many local maxima
  and infinitely many local minima. If moreover $x_n\ne x_{n+1}$ for
  all large~$n$, the local extrema are strict.
\end{lemma}

\begin{proof}
  ``Eventually non-decreasing'' is the negation of ``$d_n<0$ for
  infinitely many~$n$'', and ``eventually non-increasing'' the
  negation of ``$d_n>0$ for infinitely many~$n$''. Hence ``eventually
  monotone'' $=$ ``$\neg(d_n<0\text{ i.o.})\vee\neg(d_n>0\text{
  i.o.})$'', whose negation is exactly~(ii); and~(iii) is the same
  statement spelled out. This gives (i)$\iff$(ii)$\iff$(iii).

  Assume~(ii). Pick any $n_1$ with $d_{n_1}<0$. By~(ii) there is a
  smallest $m>n_1$ with $d_m>0$; then $d_{m-1}\le 0<d_m$, so
  $x_{m-1}\ge x_m\le x_{m+1}$, i.e.\ $m$ is a local minimum.
  By~(ii) again there is a later index with $d<0$, and the smallest
  such yields, in the same way, a local maximum beyond~$m$. Iterating
  produces infinitely many of each. Strictness under $x_n\ne x_{n+1}$
  is immediate, since then all the inequalities above are strict.
\end{proof}

By Lemma~\ref{lem:reduction}, Conjecture~\ref{conj:main} is equivalent
to the conjunction of
\begin{equation}\label{eq:two}
  \text{(A)}\quad (w_{n,p})\text{ is not eventually non-decreasing},
  \qquad
  \text{(B)}\quad (w_{n,p})\text{ is not eventually non-increasing}.
\end{equation}
Section~\ref{sec:A} proves~(A) unconditionally.
Sections~\ref{sec:transfer}--\ref{sec:status} reduce~(B) to the
oscillation of strategy~\All.

\section{Statement (A): \texorpdfstring{$w_{n,p}$}{w(n,p)} is not
  eventually non-decreasing}\label{sec:A}

\begin{proposition}\label{prop:A}
  Let $p\in(0,\tfrac12)$. There is no $M$ such that $w_{n}\ge w_{n-1}$
  for all $n>M$. That is, $w_n<w_{n-1}$ for infinitely many~$n$.
\end{proposition}

The mechanism is the following. On a stretch where $w$ is
non-decreasing, the suffix-maximum in~\eqref{eq:bellman} collapses to
its \emph{right} endpoint $w_{n-1}$, so the Bellman equation reduces to
the strategy-\One{} recursion \ref{f:strat}. Its one-step increment is
$p^n(1-w_{n-1})-q^n w_{n-1}$, in which the ``lose-at-once'' term
$q^n w_{n-1}$ (all $n$ coins tails) dominates the ``win-at-once'' term
$p^n(1-w_{n-1})$ because $q>p$. Hence the value must eventually fall.

\begin{proof}
  Suppose, for contradiction, that $w_n\ge w_{n-1}$ for all $n>M$,
  i.e.\ $w$ is non-decreasing on $\{M,M+1,\dots\}$. Fix
  $n>M+1$ and split the Bellman sum~\eqref{eq:bellman} at $j=M$:
  \[
    w_n \;=\; p^n
    \;+\; \underbrace{\sum_{j=1}^{M-1}\binom nj p^{n-j}q^j\,S_{j,n}}_{=:T_n}
    \;+\; \sum_{j=M}^{n-1}\binom nj p^{n-j}q^j\,S_{j,n},
    \qquad S_{j,n}:=\max_{j\le m\le n-1}w_m .
  \]
  For $j\ge M$ the index range $\{j,\dots,n-1\}$ lies in
  $\{M,M+1,\dots\}$, where $w$ is non-decreasing, so $S_{j,n}=w_{n-1}$.
  For every~$j$ the range contains $n-1$, hence $w_{n-1}\le
  S_{j,n}\le 1$ by~\ref{f:bounds}. Writing
  $U_n:=\sum_{j=1}^{M-1}\binom nj p^{n-j}q^j$ and using
  \ref{f:binom},
  \[
    w_n = p^n + T_n + \bigl(1-p^n-q^n-U_n\bigr)\,w_{n-1},
  \]
  and therefore
  \begin{equation}\label{eq:incrA}
    w_n-w_{n-1}
    \;=\; p^n\,(1-w_{n-1}) \;-\; q^n\,w_{n-1}
          \;+\;\bigl(T_n-U_n w_{n-1}\bigr).
  \end{equation}
  Since $w_{n-1}\le S_{j,n}\le 1$, every summand of $T_n$ lies between
  $\binom nj p^{n-j}q^j w_{n-1}$ and $\binom nj p^{n-j}q^j$, so
  $0\le T_n-U_n w_{n-1}\le U_n$. Moreover, for $1\le j\le M-1$,
  \[
    \binom nj p^{n-j}q^j \le n^{\,M-1}\,p^{\,n-M+1},
  \]
  because $\binom nj\le n^j\le n^{M-1}$, $q^j\le 1$, and
  $p^{n-j}\le p^{n-M+1}$ (smaller exponent, $0<p<1$). Hence
  \[
    U_n \;\le\; (M-1)\,n^{M-1}p^{\,n-M+1} \;=\; C_1\,n^{M-1}p^{\,n},
    \qquad C_1:=(M-1)\,p^{1-M}.
  \]
  Dropping the nonnegative term $p^n(1-w_{n-1})\le p^n$ as an upper
  bound and using $w_{n-1}\ge w_M=:c>0$ (non-decreasing, and $w_M>0$
  by~\ref{f:bounds}), \eqref{eq:incrA} gives
  \[
    w_n-w_{n-1}
    \;\le\; p^n + U_n - q^n w_{n-1}
    \;\le\; \bigl(1+C_1 n^{M-1}\bigr)\,p^n \;-\; c\,q^n .
  \]
  The ratio of the two terms is
  $\bigl(1+C_1 n^{M-1}\bigr)c^{-1}(p/q)^n\to 0$ as $n\to\infty$, since
  $p/q<1$ by~\ref{f:pq} kills the polynomial factor. Thus there is
  $n_0$ such that $w_n-w_{n-1}<0$ for all $n\ge n_0$, contradicting
  $w_n\ge w_{n-1}$ for $n>M$.
\end{proof}

\begin{remark}
  Proposition~\ref{prop:A} already refutes the natural first guess
  ``for every $p$ there is an $N$ with $w_{n,p}$ increasing for
  $n>N$''. It also shows that the main manuscript's first-order
  Corollary --- ``$w_{n,p}$ is eventually increasing'' --- is
  genuinely a \emph{first-order} statement: the exact sequence is not
  eventually increasing. The obstruction $q^n w_{n-1}$ is exponentially
  small in~$n$ but, for fixed $p<\tfrac12$, eventually beats every
  term that a finite $\delta$-expansion around $p=\tfrac12$ can see.
\end{remark}

\section{An exact transfer lemma for non-increasing
  tails}\label{sec:transfer}

We now turn to statement~(B). The key is that on a non-increasing
stretch the suffix-maximum in~\eqref{eq:bellman} collapses to its
\emph{left} endpoint, turning the Bellman equation into the
\emph{linear} strategy-\All{} recursion.

\begin{lemma}[Transfer to strategy \All]\label{lem:transfer}
  Let $p\in(0,1)$ and suppose $w_n\ge w_{n+1}$ for all $n\ge M$
  (some $M\ge 1$). Define a sequence $(x_n)_{n\ge 0}$ by
  \[
    x_0:=1,\qquad
    x_j:=\max_{j\le m\le M}w_m \quad (1\le j\le M-1),\qquad
    x_j:=w_j \quad (j\ge M).
  \]
  Then $x_n=w_n$ for all $n\ge M$, and $(x_n)$ satisfies the exact
  strategy-\All{} recursion
  \begin{equation}\label{eq:xrec}
    x_n \;=\; \sum_{j=0}^{n-1}\binom nj p^{\,n-j}q^{\,j}\,x_j
    \qquad\text{for every } n>M.
  \end{equation}
\end{lemma}

\begin{proof}
  Fix $n>M$ and consider $S_{j,n}=\max_{j\le m\le n-1}w_m$ in the
  Bellman equation~\eqref{eq:bellman}.
  \begin{itemize}[leftmargin=1.6em]
  \item For $M\le j\le n-1$: the range $\{j,\dots,n-1\}$ lies in
    $\{M,M+1,\dots\}$, where $w$ is non-increasing, so the maximum is
    at the left endpoint, $S_{j,n}=w_j=x_j$.
  \item For $1\le j\le M-1$: split
    $\{j,\dots,n-1\}=\{j,\dots,M-1\}\cup\{M,\dots,n-1\}$. On the second
    block $w$ is non-increasing, so its maximum is $w_M$; hence
    $S_{j,n}=\max\bigl(\max_{j\le m\le M-1}w_m,\;w_M\bigr)
    =\max_{j\le m\le M}w_m=x_j$, independent of~$n$.
  \end{itemize}
  Substituting into~\eqref{eq:bellman} and writing the term $p^n$ as
  $\binom n0 p^n q^0 x_0$ (with $x_0=1$) yields
  $w_n=\sum_{j=0}^{n-1}\binom nj p^{n-j}q^j x_j$. Since $x_n=w_n$ for
  $n\ge M$, this is~\eqref{eq:xrec}.
\end{proof}

\begin{remark}
  For $M=1$ the lemma specialises to the reduction lemma recorded in
  \texttt{CONVERSATION\_LOG.md} (and \texttt{localmaximum.tex}): if
  $w$ is non-increasing on all of $\N$, then $w_n=b_n$ for every~$n$.
  Lemma~\ref{lem:transfer} is the version needed when $w$ is only
  \emph{eventually} non-increasing: $w$ then equals, on its tail, a
  strategy-\All{} sequence with finitely many initial values replaced.
\end{remark}

\paragraph{Generating-function form.} Let $X(z):=\sum_{n\ge
0}x_nz^n/n!$ be the exponential generating function of the transferred
sequence. The binomial convolution underlying~\eqref{eq:xrec} has,
exactly as in \texttt{strategy\_all\_asymptotics.tex}, the EGF
$e^{pz}X(qz)$; subtracting the diagonal term $q^nx_n$ (EGF $X(qz)$)
gives that $\sum_{j=0}^{n-1}\binom nj p^{n-j}q^j x_j$ has EGF
$(e^{pz}-1)X(qz)$. Since~\eqref{eq:xrec} holds for all $n>M$ but not
necessarily for $n\le M$, the difference $X(z)-(e^{pz}-1)X(qz)$ is a
polynomial of degree at most~$M$:
\begin{equation}\label{eq:Xfunc}
  X(z) \;=\; P(z) \;+\; \bigl(e^{pz}-1\bigr)\,X(qz),
  \qquad \deg P\le M .
\end{equation}
This is precisely the $q$-difference equation~(1.3) of
\texttt{strategy\_all\_asymptotics.tex}, with the constant $1$ replaced
by the polynomial~$P$ (which encodes the finitely many ``wrong''
initial values $x_1,\dots,x_M$; for the genuine strategy-\All{} EGF
$B$ one has $P\equiv 1$). Iterating~\eqref{eq:Xfunc}, with the
remainder $\prod_{i<K}(e^{pq^iz}-1)\,X(q^Kz)\to 0$ as $K\to\infty$
because each factor $e^{pq^iz}-1\to 0$,
\begin{equation}\label{eq:Xseries}
  X(z) \;=\; \sum_{k\ge 0}\;
  \Bigl[\prod_{i=0}^{k-1}\bigl(e^{p q^i z}-1\bigr)\Bigr]\,P(q^k z).
\end{equation}

The point of~\eqref{eq:Xseries} is structural: the large-$z$ behaviour
of $X$ is driven by exactly the same product
$\prod_i(e^{pq^iz}-1)$ as that of the strategy-\All{} EGF~$B$; the
polynomial $P$ enters only as a degree-$\le M$ modulation of the
summands. In particular the saddle-point / Mellin analysis of
\texttt{strategy\_all\_asymptotics.tex} applies verbatim, and the
oscillatory poles --- located at $s_\ell=2\pi i\ell/\log(1/q)$,
$\ell\in\Z$, and produced by the geometric factor $1/(1-q^{-s})$ that
the $q$-dilation $z\mapsto qz$ generates --- are unchanged. Only the
\emph{values} of the Fourier coefficients depend on~$P$.

\section{Statement (B), conditional on the oscillation of strategy
  \All}\label{sec:status}

\begin{hypothesis}[\All-oscillation]\label{hyp:osc}
  For every $p\in(0,\tfrac12)$ the strategy-\All{} value satisfies
  $b_{n,p}=\Phi_p(\log n)\,(1+o(1))$ as $n\to\infty$, where $\Phi_p$ is
  continuous, periodic with period $\log(1/q)$, and \emph{non-constant};
  equivalently $\liminf_n b_{n,p}<\limsup_n b_{n,p}$.
\end{hypothesis}

\begin{hypothesis}[transferred \All-oscillation]\label{hyp:osc2}
  For every $p\in(0,\tfrac12)$ and every $M\ge 1$, any sequence
  $(x_n)$ produced by Lemma~\ref{lem:transfer} from an eventually
  non-increasing tail of $(w_{n,p})$ is \emph{not} eventually monotone.
\end{hypothesis}

\begin{theorem}[conditional]\label{thm:conditional}
  Assume Hypothesis~\ref{hyp:osc2}. Then for every $p\in(0,\tfrac12)$
  the sequence $(w_{n,p})$ is not eventually non-increasing. Together
  with Proposition~\ref{prop:A} this gives Conjecture~\ref{conj:main}:
  $(w_{n,p})$ is not eventually monotone, hence has infinitely many
  local maxima and minima.
\end{theorem}

\begin{proof}
  Suppose $(w_{n,p})$ were non-increasing on $\{M,M+1,\dots\}$. By
  Lemma~\ref{lem:transfer} the transferred sequence $(x_n)$ satisfies
  $x_n=w_n$ for $n\ge M$; in particular $(x_n)$ is non-increasing for
  $n\ge M$, hence \emph{eventually monotone}. This contradicts
  Hypothesis~\ref{hyp:osc2}. So $(w_{n,p})$ is not eventually
  non-increasing, which is statement~(B) of~\eqref{eq:two}. Combining
  with statement~(A) (Proposition~\ref{prop:A}) and
  Lemma~\ref{lem:reduction} proves the theorem.
\end{proof}

So the entire conjecture rests on Hypothesis~\ref{hyp:osc2}. We now
record what is known about it.

\paragraph{Status of Hypothesis~\ref{hyp:osc}.} The companion note
\texttt{strategy\_all\_asymptotics.tex} establishes:
\begin{itemize}[leftmargin=1.6em]
\item \emph{Rigorously}: the functional equation
  $B(z)=1+(e^{pz}-1)B(qz)$, the iterated series, the surjection
  formula, and the probabilistic interpretation of $b_{n,p}$.
\item \emph{By a saddle-point/Mellin derivation that is not yet a
  full proof}: the log-periodic asymptotics
  $b_{n,p}=\Phi_p(\log n)(1+o(1))$ with Fourier expansion
  $b_{n,p}=\sum_{\ell}\gamma_\ell(p)\,n^{\,2\pi i\ell/\log(1/q)}(1+o(1))$.
  Two analytic gaps remain, both ``standard in spirit''
  (Flajolet--Sedgewick): (a) Hayman-admissibility of the entire
  function $B$ --- which is of order~$1$ and type~$1$ (indeed
  $B(r)\le e^r$, since $0\le b_{n,p}\le1$), the delicate point being
  the non-constant log-periodic factor rather than the growth rate;
  (b) justification of the Mellin contour shift, i.e.\ the strip of
  analyticity.
\item \emph{Numerically}: the predicted oscillation amplitude is
  confirmed across fourteen orders of magnitude in~$p$.
\end{itemize}
A pleasant point is that, \emph{granted} the Mellin framework, the
non-constancy of $\Phi_p$ is not a separate miracle. The note computes
the Fourier coefficients as
\[
  c_\ell \;=\; -\frac{1}{\log(1/q)}\;p^{-s_\ell}\,
  \Gamma(s_\ell)\,\zeta(s_\ell+1),
  \qquad s_\ell=\frac{2\pi i\ell}{\log(1/q)} .
\]
For $\ell\ne 0$ the argument $s_\ell$ is purely imaginary and nonzero;
$\Gamma$ has no zeros, $p^{-s_\ell}\ne 0$, and
$\zeta(s_\ell+1)=\zeta(1+it)\ne 0$ because the Riemann zeta function
has no zeros on the line $\Re s=1$ (the same non-vanishing that
underlies the prime number theorem). Hence $c_\ell\ne 0$ for all
$\ell\ne 0$: \emph{within the Mellin framework, the oscillation of
$b_{n,p}$ is provably genuine.} Thus Hypothesis~\ref{hyp:osc} is
``true modulo gaps (a)--(b)''.

\paragraph{From Hypothesis~\ref{hyp:osc} to
Hypothesis~\ref{hyp:osc2}.} What
Theorem~\ref{thm:conditional} actually needs is the slightly stronger
Hypothesis~\ref{hyp:osc2}, about solutions of the \emph{$P$-modified}
equation~\eqref{eq:Xfunc}. By~\eqref{eq:Xseries} such an $X$ has the
same dominant product structure and the same oscillatory poles
$s_\ell$ as~$B$; the Mellin analysis yields a Fourier expansion
\[
  x_n \;=\; \sum_{\ell\in\Z}\gamma_\ell^X(p)\,
  n^{\,2\pi i\ell/\log(1/q)}\,(1+o(1)),
\]
with the \emph{same} period $\log(1/q)$ and with coefficients
$\gamma_\ell^X$ that differ from $\gamma_\ell$ only through a
$P$-dependent factor (the Mellin transform of the $P$-part). A
non-constant such expansion is not eventually monotone, because a
non-constant continuous periodic function attains its maximum and
minimum on every period. Hence Hypothesis~\ref{hyp:osc2} holds
\emph{provided} the modulation by~$P$ does not annihilate the
$\ell=\pm 1$ coefficient.

\section{The remaining gap, and why the two sides are
  asymmetric}\label{sec:gap}

Assembling Sections~\ref{sec:A}--\ref{sec:status}, a complete proof of
Conjecture~\ref{conj:main} is reduced to a finite, well-delimited list:
\begin{enumerate}[label=(G\arabic*),leftmargin=3em]
\item\label{g:hayman} \emph{Hayman-admissibility} of the entire
  function $B(z)$ (companion-note gap~(a)).
\item\label{g:mellin} \emph{Mellin contour justification}: the strip
  of analyticity and the contour shift (companion-note gap~(b)).
\item\label{g:P} \emph{Non-vanishing under modulation}: the
  $P$-dependent factor multiplying the $\ell=\pm 1$ pole-residue does
  not vanish for the specific polynomial $P$ that arises, via
  Lemma~\ref{lem:transfer} and~\eqref{eq:Xfunc}, from an eventually
  non-increasing tail of $(w_{n,p})$.
\end{enumerate}
Gaps \ref{g:hayman}--\ref{g:mellin} are inherited from the companion
note and are technical rather than conceptual. Gap~\ref{g:P} is the
one genuinely new analytic point introduced here. It cannot be waved
away: $P$ is not free, but it is also not the constant $1$, and one
must exclude an accidental cancellation. We expect no cancellation ---
it would be a non-generic coincidence, and $\gamma_{\pm1}$ already
carries the structurally non-zero factor $\Gamma(s_{\pm1})\zeta(1\pm
it)$ --- but a proof requires either an explicit computation of the
$P$-factor or a softer argument (for instance, showing that a
convergent solution of~\eqref{eq:xrec} would force constraints on
$x_0,\dots,x_M$ incompatible with $x_n=w_{n,p}\ge b_{n,p}$).
Section~\ref{sec:G3} carries out the first steps of that softer
argument: it removes the dependence on the specific~$P$ and reduces
gap~\ref{g:P} to a single, universal non-degeneracy.

\begin{remark}[why (A) is elementary but (B) is not]
  The asymmetry between the two halves is structural. On an
  \emph{increasing} tail the suffix-maximum collapses to the
  \emph{right} endpoint $w_{n-1}$, turning~\eqref{eq:bellman} into the
  strategy-\One{} recursion: a contractive linear recursion whose
  one-step increment has the \emph{sign-definite} dominant term
  $-q^n w_{n-1}$. That is what makes Proposition~\ref{prop:A}
  elementary. On a \emph{decreasing} tail the suffix-maximum collapses
  to the \emph{left} endpoint $w_j$, turning~\eqref{eq:bellman} into
  the strategy-\All{} recursion: a $q$-dilation-driven linear
  recursion whose solutions are log-periodic, with \emph{no}
  sign-definite drift. Ruling out monotonicity there is exactly
  ``a non-constant periodic function is not monotone'', and the
  non-constancy is the analytic content carried by
  $\zeta(1+it)\ne 0$. There is, correspondingly, no elementary
  contradiction available for case~(B): the crude estimates (weights
  concentrating at $j\approx qn$, etc.) are all consistent with a
  monotone convergent tail; only the finer Fourier structure breaks it.
\end{remark}

\begin{remark}[a weaker, more accessible statement]
  A weaker companion to Conjecture~\ref{conj:main} --- that
  $(w_{n,p})$ has \emph{at least one} strict local maximum, i.e.\ is
  not globally non-increasing --- follows from the $M=1$ case of
  Lemma~\ref{lem:transfer}: if $w$ were non-increasing on all of $\N$
  then $w_n=b_n$ for every~$n$, so it would suffice that $b_{n,p}$ is
  not monotone non-increasing. For specific $p$ this is rigorous by
  exact rational computation (e.g.\ $p=\tfrac{21}{50}$, where
  $w_9>w_8$ and $w_9>w_{10}$; see the main manuscript and
  \texttt{Manuscript/proof\_intermediate\_p042.tex}). A uniform
  statement over all $p\in(0,\tfrac12)$ again needs
  Hypothesis~\ref{hyp:osc}.
\end{remark}

\begin{remark}\label{rem:strictness}
  Lemma~\ref{lem:reduction} delivers \emph{strict} extrema once
  $w_n\ne w_{n+1}$ for all large~$n$. For $p<\tfrac12$ this is expected
  and holds in all computed data; a clean sufficient condition would
  be that $w_{n,p}$, a rational function of $p$ of growing degree,
  never has two equal consecutive values --- plausibly provable by a
  degree/positivity argument but not pursued here. Absent it,
  Conjecture~\ref{conj:main} still holds with ``strict'' replaced by
  the non-strict notion of local extremum.
\end{remark}

\section{Progress on gap~(G3): removing the dependence on
  \texorpdfstring{$P$}{P}}\label{sec:G3}

Throughout this section we work under the hypothesis to be refuted:
$(w_{n,p})$ is non-increasing on $\{M,M+1,\dots\}$. By
Lemma~\ref{lem:transfer} the transferred sequence satisfies $x_n=w_n$
for $n\ge M$; being non-increasing and bounded it converges,
\[
  x_n=w_n\;\longrightarrow\;L\qquad(n\to\infty),\qquad L\ge 0.
\]
Gap~\ref{g:P} asks us to contradict this. We show that the obstruction
can be detached from the specific polynomial~$P$.

\subsection{An elementary comparison: the sequence $g=w-b$}

\begin{lemma}\label{lem:g}
  Put $g_n:=x_n-b_n$, with $b$ the strategy-\All{} value
  of~\ref{f:strat}. Then
  \begin{enumerate}[label=\rm(\roman*),leftmargin=2.4em]
  \item $g_0=0$, and $g_n\ge 0$ for every $n\ge 0$;
  \item $g$ satisfies the homogeneous strategy-\All{} recursion
    $g_n=\sum_{j=0}^{n-1}\binom nj p^{n-j}q^j g_j$ for every $n>M$;
  \item the running maximum is non-increasing: $g_n<\max_{0\le
    j<n}g_j$ for every $n>M$, so $\max_{0\le j\le n}g_j=\max_{0\le
    j\le M}g_j=:G^\ast$ for all $n\ge M$; in particular $0\le g_n\le
    G^\ast$ for all $n$.
  \end{enumerate}
\end{lemma}

\begin{proof}
  (i) $g_0=x_0-b_0=1-1=0$. For $j\ge M$, $g_j=w_j-b_j\ge 0$ since the
  optimum dominates strategy~\All{} (\ref{f:strat}). For $1\le j\le
  M-1$, Lemma~\ref{lem:transfer} gives $x_j=\max_{j\le m\le M}w_m\ge
  w_j\ge b_j$, so $g_j\ge 0$.

  (ii) Both $x$ and $b$ satisfy
  $y_n=\sum_{j=0}^{n-1}\binom nj p^{n-j}q^j y_j$ for $n>M$ --- $x$ by
  Lemma~\ref{lem:transfer}, $b$ for every $n\ge 1$. Subtract.

  (iii) Fix $n>M$. By~(i) all $g_j\ge 0$, so by~(ii) and~\ref{f:binom},
  \[
    g_n=\sum_{j=0}^{n-1}\binom nj p^{n-j}q^j g_j
    \le\Bigl(\max_{0\le j<n}g_j\Bigr)\sum_{j=0}^{n-1}\binom nj
    p^{n-j}q^j=(1-q^n)\max_{0\le j<n}g_j<\max_{0\le j<n}g_j .
  \]
  Hence the running maximum never increases past $n=M$.
\end{proof}

\begin{proposition}\label{prop:dichotomy}
  Under the standing hypothesis, exactly one of the following holds.
  \begin{enumerate}[label=\rm(\alph*),leftmargin=2.4em]
  \item $g_n=0$ for all $n\ge M$. Then $w_n=b_n$ for all $n\ge M$;
    since $b_{n,p}$ is not eventually monotone
    (Hypothesis~\ref{hyp:osc}), neither is $(w_{n,p})$, contradicting
    the standing hypothesis.
  \item $g_n>0$ for some $n\ge M$. Then $g$ is a non-zero,
    non-negative, bounded solution of the homogeneous strategy-\All{}
    recursion (for $n>M$) with $g_0=0$.
  \end{enumerate}
  Hence gap~\ref{g:P} is equivalent to excluding case~(b).
\end{proposition}

The $g\equiv 0$ branch is thus closed rigorously and elementarily; it
uses only Hypothesis~\ref{hyp:osc}. In case~(b), $g$ would --- since
$x_n\to L$ and $b_n$ oscillates --- have to oscillate in exact
anti-phase with $b$, with asymptotic profile $L-\Phi$. Non-negativity
alone does not forbid this: by Lemma~\ref{lem:g}(iii) it only yields
$0\le g_n\le G^\ast$, and a non-negative log-periodic function is
admissible. The oscillation must be defeated through its Fourier
structure.

\subsection{Auxiliary functions $B_r$ and the elimination of $P$}

Write the modulating polynomial of~\eqref{eq:Xfunc} as
$P(z)=\sum_{r=0}^M a_r z^r$, and set
$\phi_k(z):=\prod_{i=0}^{k-1}(e^{pq^iz}-1)$, so that the iterated
series~\eqref{eq:Xseries} is $X(z)=\sum_{k\ge0}\phi_k(z)P(q^kz)$.
Expanding $P(q^kz)=\sum_r a_rq^{rk}z^r$ and exchanging the sums,
\begin{equation}\label{eq:Xdecomp}
  X(z)=\sum_{r=0}^M a_r\,z^r\,B_r(z),
  \qquad
  B_r(z):=\sum_{k\ge0}q^{rk}\phi_k(z).
\end{equation}

\begin{lemma}\label{lem:Br}
  Each $B_r$ is entire and satisfies
  $B_r(z)=1+q^{r}(e^{pz}-1)B_r(qz)$, with $B_0=B$.
\end{lemma}

\begin{proof}
  $\phi_0\equiv1$ and $\phi_k(z)=(e^{pz}-1)\phi_{k-1}(qz)$, so
  $B_r(z)=1+\sum_{k\ge1}q^{rk}(e^{pz}-1)\phi_{k-1}(qz)
  =1+q^{r}(e^{pz}-1)\sum_{k\ge0}q^{rk}\phi_k(qz)$. Convergence and
  entirety are as for $B$ in the companion note.
\end{proof}

Identity~\eqref{eq:Xdecomp} is the decisive bookkeeping step: it
expresses the transferred EGF as a fixed linear combination, with the
coefficients $a_r$ of~$P$, of the \emph{$P$-independent} functions
$z^rB_r(z)$.

\subsection{Reduction of (G3) to a $P$-independent non-degeneracy}

At the level of the saddle-point analysis of the companion note ---
hence modulo gaps~\ref{g:hayman}--\ref{g:mellin} --- each $B_r$ has a
log-periodic profile. In $B_r(z)=\sum_kq^{rk}\phi_k(z)$ the dominant
index is $k\approx K(z)=\log(pz)/\log(1/q)$, where
$q^{rK(z)}=(pz)^{-r}$ and the saddle is $O(1)$ wide; thus
\begin{equation}\label{eq:Brasymp}
  B_r(z)\;\sim\;(pz)^{-r}\,e^{z}\,\Psi_r(\log z),
\end{equation}
with $\Psi_r$ continuous, periodic of period $\log(1/q)$, and
$\Psi_0=\Phi$. Substituting~\eqref{eq:Brasymp} into~\eqref{eq:Xdecomp},
the factors $z^r$ cancel $(pz)^{-r}$ up to $p^{-r}$:
\begin{equation}\label{eq:Xprofile}
  X(z)\;\sim\;e^{z}\,\Psi^X(\log z),
  \qquad
  \Psi^X:=\sum_{r=0}^M a_r\,p^{-r}\,\Psi_r .
\end{equation}
The transferred sequence converges if and only if its profile $\Psi^X$
is constant.

\begin{proposition}\label{prop:G3red}
  Assume the asymptotics~\eqref{eq:Brasymp}. If $\Psi_0,\dots,\Psi_M$
  are linearly independent modulo the constants, then the only
  convergent solution of the transferred recursion is the zero
  sequence. As $x_0=1$, the sequence $x$ does not converge, and
  gap~\ref{g:P} is closed.
\end{proposition}

\begin{proof}
  If $\Psi^X$ is constant then, by independence, $a_rp^{-r}=0$ for
  every~$r$, i.e.\ $P\equiv0$. But $P\equiv0$ in~\eqref{eq:Xfunc}
  forces, on iterating, $X(z)=\phi_K(z)X(q^Kz)\to0$, so $X\equiv0$ and
  $x\equiv0$, against $x_0=1$. Hence $\Psi^X$ is non-constant, and a
  non-constant continuous periodic function is not eventually
  monotone.
\end{proof}

The hypothesis of Proposition~\ref{prop:G3red} is genuinely
$P$-independent: it speaks only of the universal family
$(\Psi_r)_{r\ge0}$. That this is the \emph{right} reformulation, with
no loss, is confirmed by an exact observation. If $\sum_{r=0}^M
c_r\Psi_r$ were constant, set $Y(z):=\sum_{r=0}^M c_rz^rB_r(z)$;
using $B_r(z)=1+q^r(e^{pz}-1)B_r(qz)$ and $z^r=q^{-r}(qz)^r$,
\[
  Y(z)=\widetilde C(z)+(e^{pz}-1)\,Y(qz),
  \qquad \widetilde C(z):=\sum_{r=0}^M c_rz^r,
\]
the \emph{same} $q$-difference equation~\eqref{eq:Xfunc}. Hence
gap~\ref{g:P}, for every~$M$, is equivalent to
\[
  (\Psi_r)_{r\ge0}\ \text{linearly independent modulo constants}
  \iff
  \text{no non-zero polynomial }\widetilde C\text{ yields a convergent
  \All-type sequence.}
\]

\subsection{Status of (G3) after this reduction}

\begin{itemize}[leftmargin=1.6em]
\item \textbf{Gained.} Gap~\ref{g:P} no longer refers to the
  polynomial~$P$ of the specific tail of~$w$. By
  Proposition~\ref{prop:G3red} it is reduced to one universal
  non-degeneracy: linear independence, modulo constants, of the
  profiles $\Psi_r$ of the auxiliary functions
  $B_r=1+q^r(e^{pz}-1)B_r(qz)$. The $g\equiv0$ branch
  (Proposition~\ref{prop:dichotomy}(a)) is moreover fully rigorous.
\item \textbf{Still open.} The independence of the $\Psi_r$. Via the
  Mellin formalism it is the assertion that the Fourier-coefficient
  matrix $[\widehat\Psi_{r,\ell}]_{r\ge0,\,\ell\in\Z}$ has full row
  rank. We expect this --- for $r=0$ the relevant coefficients already
  carry the structurally non-zero factor $\Gamma(s_\ell)\zeta(s_\ell+1)$
  --- but it is not proved, and it inherits
  gaps~\ref{g:hayman}--\ref{g:mellin}.
\item \textbf{What does not suffice.} Non-negativity of $g$
  (Lemma~\ref{lem:g}): the running-maximum bound gives only $0\le
  g_n\le G^\ast$, compatible with a non-negative oscillation of
  profile $L-\Phi$.
\end{itemize}
In short, gap~\ref{g:P} is now sharp and $P$-free: it asserts that no
non-trivial polynomial modulation of the strategy-\All{}
$q$-difference equation produces a convergent sequence.

\section{A renormalisation route to Hypothesis~\ref{hyp:osc2}}
\label{sec:renorm}

Sections~\ref{sec:status}--\ref{sec:G3} reduce
Hypothesis~\ref{hyp:osc2} to the Mellin analysis of the companion
note. This section sketches an alternative, more self-contained route:
a scale-by-scale induction that replaces the Mellin/Hayman machinery
(gaps~\ref{g:hayman}--\ref{g:mellin}) by a single elementary
\emph{propagation estimate}. It does not remove the
non-degeneracy~\ref{g:P}, but it isolates it as one finite,
numerically accessible base case.

Throughout, $(x_n)$ denotes a bounded solution of the strategy-\All{}
recursion for $n>M$ --- a transferred sequence of
Lemma~\ref{lem:transfer}, or $b$ itself.

\subsection{The recursion as a smoothing in log-scale}

\begin{lemma}\label{lem:smooth}
  For $n>M$, with $J_n\sim\mathrm{Bin}(n,q)$,
  \[
    x_n=\frac{E[x_{J_n}]}{1+q^n}.
  \]
\end{lemma}

\begin{proof}
  Since $\binom nj p^{n-j}q^j=P(J_n=j)$,
  $\sum_{j=0}^{n-1}\binom nj p^{n-j}q^j x_j=E[x_{J_n}]-P(J_n=n)x_n
   =E[x_{J_n}]-q^nx_n$. Equate with $x_n$ and solve.
\end{proof}

As $E[J_n]=qn$ and $\operatorname{Var}(J_n)=npq$, the index $J_n$
concentrates at $qn$ with relative spread $\sqrt{p/(qn)}\to0$. Pass to
the logarithmic variable: put
\[
  t:=\log n,\qquad L:=\log(1/q),\qquad \omega:=2\pi/L,
\]
and let $y$ be the log-linear interpolation of $n\mapsto x_n$, so
$y(\log n)=x_n$. The decisive arithmetic fact is
\[
  \log(qn)=\log n-L:
\]
one application of the recursion shifts the log-scale down by
\emph{exactly one period}. By the delta method $\log J_n$ has mean
$t-L+O(1/n)$ and standard deviation
$\sigma(t)=\sqrt{p/q}\;e^{-t/2}\,(1+o(1))\to0$. Hence
Lemma~\ref{lem:smooth} reads, up to the negligible factor
$1/(1+q^n)=1-O(e^{-Le^{t}})$,
\begin{equation}\label{eq:ysmooth}
  y(t)=E\bigl[y\bigl(t-L+\sigma(t)\,Z_t\bigr)\bigr]+\varepsilon(t),
\end{equation}
where $Z_t$ is the standardised $\log J_n$ and $\varepsilon(t)$
collects the $1/(1+q^n)$ factor and the indices $j\le M$ (total
binomial weight $O(n^Mp^n)$), so $|\varepsilon(t)|=O(e^{-c\,e^{t}})$.
In words: \emph{$y$ at log-scale $t$ equals $y$ at log-scale $t-L$,
smoothed by a kernel of width $\sigma(t)\to0$.}

\subsection{The local fundamental coefficient}

Fix a window $\chi\in C_c^\infty(\mathbb R)$ with
\[
  \int_{\mathbb R}\chi=1,\qquad
  \hat\chi(2\pi):=\int_{\mathbb R}\chi(u)\,e^{-2\pi i u}\,du=0
\]
(one linear constraint, easily satisfied). Define the \emph{local
fundamental coefficient} of $y$ at log-scale $t$:
\begin{equation}\label{eq:Adef}
  A(t):=\frac1L\int_{\mathbb R} y(s)\,e^{-i\omega s}\,
        \chi\!\Bigl(\tfrac{s-t}{L}\Bigr)\,ds .
\end{equation}

\begin{lemma}\label{lem:Aconv}
  If $x_n\to x_\infty$, then $A(t)\to0$ as $t\to\infty$.
\end{lemma}

\begin{proof}
  For large $t$ the window sees only $s$ with $y(s)=x_\infty+o(1)$.
  The leading term is
  $\frac{x_\infty}{L}\int e^{-i\omega s}\chi(\tfrac{s-t}L)\,ds
   =x_\infty e^{-i\omega t}\hat\chi(\omega L)
   =x_\infty e^{-i\omega t}\hat\chi(2\pi)=0$ since $\omega L=2\pi$;
  the $o(1)$ part contributes $o(1)$.
\end{proof}

Thus a non-vanishing $A$ certifies that $(x_n)$ does \emph{not}
converge, hence --- being bounded --- is not eventually monotone.

\subsection{The propagation estimate}

\begin{proposition}[Propagation]\label{prop:prop}
  There exist $\kappa(t)\in\mathbb C$ and a remainder $r(t)$ with
  \begin{equation}\label{eq:prop}
    A(t+L)=\kappa(t)\,A(t)+r(t),
    \qquad |\kappa(t)-1|=O(e^{-t}),\qquad |r(t)|=O(e^{-t/2}),
  \end{equation}
  the implied constants depending only on $p$, $\chi$ and
  $\sup_n|x_n|$.
\end{proposition}

\begin{proof}[Proof sketch]
  Insert~\eqref{eq:ysmooth} into~\eqref{eq:Adef} evaluated at $t+L$.
  The substitution $s\mapsto s+L$ turns the period shift into the
  identity on the phase, because $e^{-i\omega L}=e^{-2\pi i}=1$;
  removing the smoothing variable $s\mapsto s+\sigma(t)Z_t$ then
  produces the multiplier
  \[
    \kappa(t)=E\bigl[e^{\,i\omega\sigma(t)Z_t}\bigr]
            =1-\tfrac12\omega^2\sigma(t)^2+O(\sigma(t)^3)
            =1-O(e^{-t}),
  \]
  the characteristic function of the smoothing kernel at the
  fundamental frequency. The remainder $r(t)$ collects: (i) the
  displacement of the window argument by $\sigma(t)Z_t$, of order
  $O(\sigma(t))=O(e^{-t/2})$; (ii) the Edgeworth correction to the
  Gaussian approximation of $\log J_n$, $O(\sigma(t)^3)$; (iii) the
  interpolation and Euler--Maclaurin errors, $O(1/n)=O(e^{-t})$;
  (iv) $\varepsilon(t)$, super-exponentially small. The dominant term
  is~(i), giving $|r(t)|=O(e^{-t/2})$. Quantitative forms of
  (i)--(iv) are standard (Berry--Esseen / Edgeworth bounds for the
  binomial, uniform in the relevant range of $t$); the constants are
  not tracked here.
\end{proof}

The two exponents in~\eqref{eq:prop} are what matters: both
$\sum_k|\kappa(t_0+kL)-1|$ and $\sum_k|r(t_0+kL)|$ converge, being
geometric series in $e^{-L}<1$.

The factor $\kappa$ admits a fully rigorous treatment; only the
remainder $r(t)$ rests on the sketch. With $t=\log n$, the recursion
of Lemma~\ref{lem:smooth} shifts the log-scale down by exactly one
period to $\log(qn)$, and $\kappa$ is the resulting
characteristic-function-type average over the binomial law.

\begin{lemma}[the propagation factor]\label{lem:kappa}
  Let $0<q<1$, $p=1-q$, $\omega\in\mathbb R$, and
  $J_n\sim\mathrm{Bin}(n,q)$. Then
  \[
    \kappa(n):=E\bigl[(J_n/(qn))^{\,i\omega}\,
                      \mathbf 1_{\{J_n\ge1\}}\bigr]
  \]
  satisfies, for every $n\ge1$,
  \[
    \kappa(n)=1-\frac{p}{2qn}\,(\omega^{2}+i\omega)+R_n,
    \qquad |R_n|\le \frac{C(\omega,q)}{n^{3/2}},
  \]
  with $C(\omega,q)$ explicit. In particular $|\kappa(n)-1|=O(1/n)$,
  and $n\,|\kappa(n)-1|\to\frac{p}{2q}\,|\omega^{2}+i\omega|
  =\frac{p\,\omega\sqrt{\omega^{2}+1}}{2q}$.
\end{lemma}

\begin{proof}
  Put $V_n:=J_n/(qn)-1$, so $E[V_n]=0$ and $-1\le V_n\le1/q$. The
  central moments of $\mathrm{Bin}(n,q)$ give the exact values
  $E[V_n^{2}]=\operatorname{Var}(J_n)/(qn)^{2}=p/(qn)$ and
  $E[V_n^{4}]=3p^{2}/(q^{2}n^{2})+p(1-6pq)/(q^{3}n^{3})$; hence, for
  $n\ge1/(pq)$, $E[V_n^{4}]\le4p^{2}/(q^{2}n^{2})$ and, by
  Cauchy--Schwarz,
  \begin{equation}\label{eq:V3bound}
    E[|V_n|^{3}]\le\bigl(E[V_n^{2}]\,E[V_n^{4}]\bigr)^{1/2}
                \le2\,(p/q)^{3/2}\,n^{-3/2}.
  \end{equation}

  \emph{Taylor expansion.} The map $g(v):=(1+v)^{i\omega}
  =\exp(i\omega\log(1+v))$ has $g(0)=1$, $g'(0)=i\omega$,
  $g''(0)=i\omega(i\omega-1)=-(\omega^{2}+i\omega)$, and
  $|g'''(v)|=|\omega|\sqrt{\omega^{2}+1}\,\sqrt{\omega^{2}+4}\,
  (1+v)^{-3}$. For $|v|\le\tfrac12$ one has $(1+v)^{-3}\le8$, so
  Taylor's theorem with integral remainder gives
  \begin{equation}\label{eq:gtaylor}
    \bigl|(1+v)^{i\omega}-1-i\omega v
          +\tfrac12(\omega^{2}+i\omega)v^{2}\bigr|
    \le\tfrac{M_\omega}{6}|v|^{3},
    \quad M_\omega:=8|\omega|\sqrt{\omega^{2}+1}\,\sqrt{\omega^{2}+4}.
  \end{equation}

  \emph{The tail event.} By Hoeffding's inequality
  $P(|V_n|>\tfrac12)=P(|J_n-qn|>qn/2)\le2e^{-q^{2}n/2}$. On
  $\{|V_n|\le\tfrac12\}$ one has $J_n\ge qn/2\ge1$ once $n\ge2/q$;
  on the complement $|(1+V_n)^{i\omega}|=1$ and $|V_n|\le1/q$.

  \emph{Assembling} (for $n\ge\max(2/q,1/(pq))$). Split
  $\kappa(n)=E[(1+V_n)^{i\omega}\mathbf 1_{|V_n|\le1/2}]
  +E[(1+V_n)^{i\omega}\mathbf 1_{|V_n|>1/2,\,J_n\ge1}]$; the second
  term has modulus $\le P(|V_n|>\tfrac12)\le2e^{-q^{2}n/2}$. Inserting
  \eqref{eq:gtaylor} into the first term,
  \[
    E[(1+V_n)^{i\omega}\mathbf 1_{|V_n|\le1/2}]
    =P(|V_n|\le\tfrac12)+i\omega\,E[V_n\mathbf 1_{|V_n|\le1/2}]
     -\tfrac12(\omega^{2}+i\omega)\,E[V_n^{2}\mathbf 1_{|V_n|\le1/2}]
     +\mathcal E_n,
  \]
  with $|\mathcal E_n|\le\tfrac{M_\omega}{6}E[|V_n|^{3}]$. Each
  truncated moment differs from the full moment ($E[V_n]=0$,
  $E[V_n^{2}]=p/(qn)$) by at most $q^{-2}P(|V_n|>\tfrac12)
  \le2q^{-2}e^{-q^{2}n/2}$, since $|V_n|\le1/q$. Collecting the
  exponentially small terms and using~\eqref{eq:V3bound},
  \[
    \kappa(n)=1-\tfrac{p}{2qn}(\omega^{2}+i\omega)+R_n,
    \qquad
    |R_n|\le\tfrac{M_\omega}{3}(p/q)^{3/2}n^{-3/2}
            +c(\omega,q)\,e^{-q^{2}n/2}.
  \]
  As $e^{-q^{2}n/2}=O(n^{-3/2})$ this is the claim; the finitely many
  $n<\max(2/q,1/(pq))$ are absorbed by enlarging $C(\omega,q)$. The
  limit statement follows since $|R_n|=o(1/n)$.
\end{proof}

\begin{remark}
  Lemma~\ref{lem:kappa} makes the $\kappa$-half of
  Proposition~\ref{prop:prop} rigorous, and sharper than the sketch:
  beyond $|\kappa-1|=O(e^{-t})$ it gives the exact leading constant,
  including the imaginary part $\tfrac{p\omega}{2qn}$ produced by the
  convexity of $\log$. Only the remainder $r(t)$ --- the
  windowed-coefficient bookkeeping --- still rests on the proof
  sketch.
\end{remark}

\subsection{The induction}

\begin{theorem}[renormalisation form]\label{thm:renorm}
  Let $(x_n)$ be a bounded solution of the strategy-\All{} recursion
  for $n>M$, and assume the propagation estimate~\eqref{eq:prop}.
  Then for each fixed $t_0$ the sequence $A(t_0+kL)$ converges as
  $k\to\infty$ to a limit $A_\infty$. Moreover
  \begin{enumerate}[label=\rm(\roman*),leftmargin=2.4em]
  \item if $x_n$ converges, then $A_\infty=0$;
  \item if $A_\infty\ne0$, then $(x_n)$ is not eventually monotone.
  \end{enumerate}
\end{theorem}

\begin{proof}
  Write $A_k:=A(t_0+kL)$, $\kappa_k:=\kappa(t_0+kL)$,
  $r_k:=r(t_0+kL)$. By~\eqref{eq:prop}, $\sum_k|\kappa_k-1|<\infty$
  and $\sum_k|r_k|<\infty$, so the products
  $\prod_{j}\kappa_j$ converge with non-zero value. Solving the
  linear recursion $A_{k+1}=\kappa_kA_k+r_k$,
  \[
    A_k=\Bigl(\prod_{j<k}\kappa_j\Bigr)A_0
       +\sum_{j<k}\Bigl(\prod_{j<i<k}\kappa_i\Bigr)r_j
       \;\xrightarrow[k\to\infty]{}\;
       A_\infty:=\Bigl(\prod_{j\ge0}\kappa_j\Bigr)A_0
       +\sum_{j\ge0}\Bigl(\prod_{i>j}\kappa_i\Bigr)r_j,
  \]
  the series converging absolutely; this proves convergence.
  (i) If $x_n\to x_\infty$ then $A(t)\to0$ by
  Lemma~\ref{lem:Aconv}, so $A_\infty=0$. (ii) If $A_\infty\ne0$ then
  $A_k\not\to0$, so by~(i) $x_n$ does not converge; a bounded
  non-convergent sequence is not eventually monotone.
\end{proof}

So, granted the propagation estimate, Hypothesis~\ref{hyp:osc2} for a
given transferred sequence $x$ is \emph{equivalent} to the single
non-degeneracy $A_\infty(x)\ne0$.

\subsection{What the route achieves}

\begin{itemize}[leftmargin=1.6em]
\item \textbf{Replaces \ref{g:hayman}--\ref{g:mellin}.} The Mellin
  contour shift and Hayman-admissibility are replaced by the single
  Proposition~\ref{prop:prop}, whose error terms (i)--(iv) are
  explicit and elementary; they are summable along the geometric
  scale precisely because $L=\log(1/q)>0$. The Mellin transform
  diagonalises the $q$-dilation in one stroke; the induction does the
  same diagonalisation approximately, scale by scale.
\item \textbf{Isolates \ref{g:P}.} By linearity $A_\infty$ depends
  linearly on the data; for a transferred $x=b+g$ (notation of
  \S\ref{sec:G3}) one has $A_\infty(x)=A_\infty(b)+A_\infty(g)$.
  Here $A_\infty(b)\ne0$ is the genuine strategy-\All{} oscillation
  --- it coincides with the fundamental Mellin coefficient
  $\gamma_1(p)$, non-zero by $\Gamma(s_1)\zeta(s_1+1)\ne0$ --- and
  ruling out a cancellation by $A_\infty(g)$ is exactly
  gap~\ref{g:P}.
\item \textbf{Makes the base case finite and computable.} Unlike the
  asymptotic gaps \ref{g:hayman}--\ref{g:mellin}, the surviving
  quantity $A_\infty$ is the limit of the \emph{explicitly
  computable} sequence $A(t_0+kL)$ of~\eqref{eq:Adef}; it is
  therefore numerically accessible. For the genuine $b$
  (Hypothesis~\ref{hyp:osc}) the route is complete modulo
  Proposition~\ref{prop:prop} and the single inequality
  $A_\infty(b)\ne0$. This picture is borne out numerically: the script
  \texttt{simulation/renorm\_validation.py} computes, for
  $p\in\{0.35,0.40,0.42,0.45\}$, both the windowed coefficient
  $A(t_0+kL)$ --- found to converge to a non-zero limit --- and the
  factor $\kappa(n)$, for which $n\,|\kappa(n)-1|$ approaches the
  constant $\tfrac{p}{2q}|\omega^{2}+i\omega|$ of
  Lemma~\ref{lem:kappa}.
\end{itemize}

\section{Summary}\label{sec:summary}

\begin{itemize}[leftmargin=1.6em]
\item \textbf{Rigorous (this note).}
  Conjecture~\ref{conj:main} is equivalent to the conjunction of~(A)
  and~(B) in~\eqref{eq:two} (Lemma~\ref{lem:reduction}).
  Statement~(A), ``$w_{n,p}$ is not eventually non-decreasing'', is
  proved unconditionally (Proposition~\ref{prop:A}). The transfer
  lemma (Lemma~\ref{lem:transfer}) is exact: on an eventually
  non-increasing tail, $w$ \emph{is} a strategy-\All{} sequence with
  finitely many initial values modified.
\item \textbf{Conditional.} Statement~(B), and hence the full
  conjecture, follows from Hypothesis~\ref{hyp:osc2}
  (Theorem~\ref{thm:conditional}) --- the log-periodic oscillation of
  the $P$-modified strategy-\All{} recursion.
\item \textbf{Open.} Hypothesis~\ref{hyp:osc2} reduces to the gaps of
  Section~\ref{sec:gap}: the two inherited, technical Mellin/Hayman
  gaps \ref{g:hayman}--\ref{g:mellin}, and gap~\ref{g:P}. The latter
  is sharpened in Section~\ref{sec:G3}: its $g\equiv0$ branch is closed
  rigorously (Proposition~\ref{prop:dichotomy}(a)), and the remainder
  is reduced to a single $P$-independent non-degeneracy --- linear
  independence, modulo constants, of the profiles of the auxiliary
  functions $B_r$ (Proposition~\ref{prop:G3red}). None is a deep
  obstruction; together they are a concrete finishing project.
\item \textbf{Alternative route.} Section~\ref{sec:renorm} recasts
  Hypothesis~\ref{hyp:osc2} as a scale-by-scale induction: it replaces
  the Mellin gaps \ref{g:hayman}--\ref{g:mellin} by one elementary
  propagation estimate (Proposition~\ref{prop:prop}) with summable
  errors, and reduces the conjecture to the single, numerically
  accessible non-degeneracy $A_\infty\ne0$.
\end{itemize}

\paragraph{Suggested use in the main manuscript.} Section~6 of the
main paper could carry, honestly:
\begin{itemize}[leftmargin=1.6em]
\item a \emph{Proposition} --- ``for $p<\tfrac12$, $w_{n,p}$ is not
  eventually monotone increasing'' --- with the short proof of
  Proposition~\ref{prop:A}; and
\item a \emph{Conjecture} --- Conjecture~\ref{conj:main} --- with the
  transfer lemma as the indication of proof and a pointer to
  \texttt{strategy\_all\_asymptotics.tex} for the analytic input.
\end{itemize}
This pairs a clean new theorem with a sharply-delimited open problem,
and corrects the impression, left by the first-order Corollary of
Section~4, that $w_{n,p}$ is eventually increasing.

\end{document}
